% from ega4-11-fr.tex
\oldpage[IV-3]{155}

précédemment permet de remplacer $A'$ par $B'$~; on peut donc supposer l'anneau
local $A'$ intégralement clos, $A$ étant un sous-anneau local de $A'$ dominé par
$A'$. Soit $A_1$ la clôture intégrale de $A$~; il est clair que
$A\subset A_1\subset A'$, et par hypothèse $A_1$ est un anneau local~; si
$\mathfrak m$, $\mathfrak m_1$, $\mathfrak m'$ sont les idéaux maximaux de
$A$, $A_1$, $A'$, on a $\mathfrak m\subset\mathfrak m_1\subset\mathfrak m'$~;
en effet, $\mathfrak m_1$ est le \emph{seul} idéal premier de $A_1$ au-dessus de
$\mathfrak m$, puisque $A_1$ est un anneau local (Bourbaki,
\emph{Alg. comm.}, chap.~V, \S~2, n$^\circ$~1, prop.~1)~; comme
$\mathfrak m'\cap A=\mathfrak m$, on a
$\mathfrak m'\cap A_1\cap A=\mathfrak m$, donc
$\mathfrak m'\cap A_1=\mathfrak m_1$. Posons
$Y_1=\operatorname{Spec}(A_1)$, $X_1=X\times_Y Y_1$, $f_1=f_{(Y_1)}$,
$\mathcal F_1=\mathcal F\otimes_Y Y_1$, et soit $x_1$ la projection de $x'$ dans
$X_1$~; désignons d'autre part par $y_1$ l'unique point fermé de $Y_1$, de sorte
que $y_1=f_1(x_1)$. Par hypothèse le morphisme
$\operatorname{Spec}(k(y_1))\to\operatorname{Spec}(k(y))$ est radiciel, d'où
l'on conclut, par la transitivité des fibres (\textbf I, 3.6.4) et
(\textbf I, 3.5.7), que le morphisme $f_1^{-1}(y_1)\to f^{-1}(y)$ est radiciel,
et en particulier que $x_1$ est le \emph{seul} point de $X_1$ dont les
projections dans $X$ et $Y_1$ soient $x$ et $y_1$ respectivement~; en outre, on
a vu que $y_1$ est le seul point de $Y_1$ dont la projection dans $Y$ soit $y$,
donc $x_1$ est le seul point de $X_1$ dont la projection dans $X$ soit $x$. Si
l'on prouve que $\mathcal F_1$ est $f_1$-plat au point $x_1$, on pourra donc
appliquer (11.5.5), d'où résultera la conclusion. On est donc ramené au cas où
$A$ lui aussi est \emph{intégralement clos}.

\medskip
\noindent\textbf{II)} \emph{Réduction au cas où $A$ et $A'$ sont des anneaux
locaux de $\mathbf Z$-algèbres de type fini intégralement clos.}

On peut considérer $A'$ comme limite inductive filtrante de ses
sous-$\mathbf Z$-algèbres (intègres) de type fini $B'_{\lambda}$~; en outre,
comme $A'$ est intégralement clos et que la clôture intégrale d'une
$\mathbf Z$-algèbre de type fini est aussi une $\mathbf Z$-algèbre de type fini
(7.8.3), on voit que $A'$ est limite inductive de ses sous-$\mathbf Z$-algèbres
de type fini $B'_{\lambda}$ \emph{intégralement closes}~; si
$\mathfrak m'_{\lambda}=\mathfrak m'\cap B'_{\lambda}$, $A'$ est aussi limite
inductive des sous-anneaux locaux
$(B'_{\lambda})_{\mathfrak m'_{\lambda}}=A'_{\lambda}$ dominés par $A'$
(5.13.3). Pour tout $\lambda$, $B'_{\lambda}\cap A$ est aussi limite inductive
de ses sous-$\mathbf Z$-algèbres de type fini $B_{\alpha\lambda}$, donc
$A=\varinjlim_{\alpha,\lambda}B_{\alpha\lambda}$, et comme plus haut on peut
remplacer $B_{\alpha\lambda}$ dans cette formule par sa clôture intégrale
(contenue par hypothèse dans $B'_{\lambda}\cap A$), puis par l'anneau local
$A_{\alpha\lambda}=(B_{\alpha\lambda})_{\mathfrak m_{\alpha\lambda}}$, où
$\mathfrak m_{\alpha\lambda}=\mathfrak m\cap B_{\alpha\lambda}
=\mathfrak m'_{\lambda}\cap B_{\alpha\lambda}$, de sorte que
$A_{\alpha\lambda}$ est dominé par $A'_{\lambda}$ et par $A$. Posons
$Y_{\alpha\lambda}=\operatorname{Spec}(A_{\alpha\lambda})$~; il résulte de
(8.9.1) qu'il existe un couple $(\alpha,\lambda)$ assez grand, un morphisme
$f_{\alpha\lambda}:X_{\alpha\lambda}\to Y_{\alpha\lambda}$ de type fini et un
$\mathcal O_{X_{\alpha\lambda}}$-Module cohérent $\mathcal F_{\alpha\lambda}$
tels que
\[
X=X_{\alpha\lambda}\times_{Y_{\alpha\lambda}}Y,
\qquad f=f_{\alpha\lambda}\times1_Y,
\qquad \mathcal F=\mathcal F_{\alpha\lambda}\otimes_{Y_{\alpha\lambda}}Y~;
\]
si $x_{\alpha\lambda}$ est la projection de $x$ dans $X_{\alpha\lambda}$, il
suffira de montrer que $\mathcal F_{\alpha\lambda}$ est
$f_{\alpha\lambda}$-plat au point $x_{\alpha\lambda}$. Comme $A'$ est limite
inductive des $A'_{\mu}$ pour $\mu\geq\lambda$, $X'$ est limite projective des
$X_{\alpha\lambda}\times_{Y_{\alpha\lambda}}Y'_{\mu}=X'_{\mu}$, et on a aussi
$\mathcal F'=\mathcal F'_{\alpha\mu}\otimes_{Y'_{\mu}}Y'$, où
$\mathcal F'_{\alpha\mu}
=\mathcal F_{\alpha\lambda}\otimes_{Y_{\alpha\lambda}}Y'_{\mu}$.
Appliquant (11.2.6), on voit qu'on peut prendre $\mu$ assez grand pour que
$\mathcal F'_{\alpha\mu}$ soit $Y'_{\mu}$-plat au point $x'_{\mu}$, projection
de $x'$ dans $X'_{\mu}$, et en outre, par construction des $A'_{\mu}$, la
projection de $x'_{\mu}$ dans $Y'_{\mu}$ est le point fermé de $Y'_{\mu}$.

\medskip
\noindent\textbf{III)} \emph{Réduction au cas où le corps résiduel $k'$ de
$A'$ est une extension finie et radicielle du corps résiduel $k$ de $A$.}

On peut en premier lieu répéter le raisonnement de la partie I) de la
démonstration de (11.4.11), tenant compte du fait que $\mathbf Z$ est un anneau
de
% IV3_B23_P155_END

